\clearpage
\begin{proofbox}
\textbf{统一设定：}设 $P=(x_0,y_0)$ 在 $H$ 上；两渐近线分别为
$\ell_1:y=(b/a)x$ 与 $\ell_2:y=-(b/a)x$，相交于 $O=(0,0)$。

\medskip
\textbf{方法一｜最快证明（约三行：标准切线＋两个交点＋行列式）。}
\begin{center}
\begin{tikzpicture}[x=1.12cm,y=1.12cm,>=stealth,every node/.style={font=\small}]
\coordinate (O) at (0,0);
\coordinate (A) at (3.24,1.98);
\coordinate (B) at (1,-0.611111);
\coordinate (P) at (2.12,0.684444);
\fill[red!7] (O)--(A)--(B)--cycle;
\draw[gray!65,->] (-0.65,0)--(4.15,0) node[right] {$x$};
\draw[gray!65,->] (0,-1.35)--(0,2.6) node[above] {$y$};
\draw[blue!65!black,densely dashed] (-0.65,-0.3972)--(3.8,2.3222);
\draw[blue!65!black,densely dashed] (-0.65,0.3972)--(2.05,-1.2528);
\draw[teal!70!black,thick,domain=-0.9:1.22,samples=90,variable=\t]
 plot ({1.8/2*(exp(\t)+exp(-\t))},{1.1/2*(exp(\t)-exp(-\t))});
\draw[red!75!black,very thick] ($(B)!-0.15!(A)$)--($(A)!-0.12!(B)$);
\draw[orange!90!black,thick] (A)--(O)--(B);
\fill (O) circle(1.5pt) node[below left=5pt] {$O$};
\fill (A) circle(1.5pt) node[above left] {$A$};
\fill (B) circle(1.5pt) node[below=9pt] {$B$};
\fill (P) circle(1.5pt) node[left=4pt] {$P$};

\draw[gray,densely dotted] (A)--(3.24,0);
\draw[gray,densely dotted] (B)--(1,0);
\node[anchor=west] at (4.55,2.05) {$A=\left(\frac a{u-v},\frac b{u-v}\right)$};
\node[anchor=west] at (4.55,0.85) {$B=\left(\frac a{u+v},-\frac b{u+v}\right)$};
\node[anchor=west] at (4.55,-0.35) {$2S=|x_Ay_B-x_By_A|$};
\end{tikzpicture}

\end{center}

令 $u=x_0/a,\ v=y_0/b$，则 $u^2-v^2=1$。
切线是 $u(x/a)-v(y/b)=1$，分别代入 $y=\pm(b/a)x$，
得到图中标注的 $A,B$ 坐标。
因为 $(u-v)(u+v)=1$，所以
\[
\boxed{S_{\triangle AOB}=\frac12|x_Ay_B-x_By_A|
=\frac{ab}{|u^2-v^2|}=ab}.
\]
\textbf{优势：}切点在左支、右支或顶点都不用分类；唯一触发信号是
\emph{切线＋两条渐近线＋面积}。

\medskip
\end{proofbox}
\clearpage
\begin{proofbox}
\textbf{方法二｜最低门槛（先看顶点，再看``一长一短''）。}
\begin{center}
% Same a,b, coordinate ranges and scale in all three panels.
\begin{tikzpicture}[x=0.8cm,y=0.8cm,>=stealth,every node/.style={font=\small}]
\foreach \rr/\shift/\caption in {1/0/{$r=1$},1.4/5.1/{$r=1.4$},1.8/10.2/{$r=1.8$}} {
\begin{scope}[shift={(\shift,0)}]
\coordinate (O) at (0,0);
\coordinate (A) at ({1.8*\rr},{1.1*\rr});
\coordinate (B) at ({1.8/\rr},{-1.1/\rr});
\coordinate (P) at ({.9*(\rr+1/\rr)},{.55*(\rr-1/\rr)});
\fill[red!7] (O)--(A)--(B)--cycle;
\draw[gray!65,->] (-.4,0)--(4,0);
\draw[gray!65,->] (0,-1.65)--(0,2.65);
\draw[blue!65!black,densely dashed] (-.3,-.1833)--(3.7,2.2611);
\draw[blue!65!black,densely dashed] (-.3,.1833)--(2.5,-1.5278);
\draw[teal!70!black,thick,domain=-.9:1.15,samples=70,variable=\t]
 plot ({.9*(exp(\t)+exp(-\t))},{.55*(exp(\t)-exp(-\t))});
\draw[orange!90!black,thick] (A)--(O)--(B);
\draw[red!75!black,very thick] ($(B)!-.12!(A)$)--($(A)!-.08!(B)$);
\fill (O) circle(1.5pt) node[below left=5pt] {$O$};
\fill (A) circle(1.5pt) node[above left] {$A$};
\fill (B) circle(1.5pt) node[below=9pt] {$B$};
\fill (P) circle(1.5pt) node[left] {$P$};
\node at (1.7,3.45) {\caption};
\node at (1.7,-2.85) {$S=ab$};
\end{scope}}
\end{tikzpicture}

\end{center}

顶点 $P=(a,0)$ 时，$AB=2b$，原点到 $AB$ 距离 $a$，故 $S=ab$。
对一般切线，先利用切线截两渐近线的\emph{中点性质} $P=(A+B)/2$。
可把两个交点写为
\[
A=(ar,br),\qquad B=(as,-bs)\qquad(r,s\ne0).
\]
于是 $P=(a(r+s)/2,\ b(r-s)/2)$；因 $P\in H$，
\[
1=\frac{(r+s)^2}{4}-\frac{(r-s)^2}{4}=rs.
\]
因此 $A=(ar,br),\ B=(a/r,-b/r)$。图示右支的正参数；
一般长度分别按 $|r|$ 与 $1/|r|$ 倍变化，进而
\[
S_{\triangle AOB}=
\frac12\left|ar\cdot\left(-\frac br\right)
-\frac ar\cdot br\right|=ab.
\]
\textbf{中点性质不能只凭图猜：}用方法一中两个交点坐标，
结合 $1/(u-v)=u+v$、$1/(u+v)=u-v$，直接得到
\[
\frac{A+B}{2}
=\left(\frac a2[(u+v)+(u-v)],\frac b2[(u+v)-(u-v)]\right)
=(au,bv)=P.
\]
\textbf{优势：}理解形状变化的原因，而不仅是算出一个数；只需``平方差＝乘积''。

\medskip
\textbf{顺带推出长度乘积：}按方法二，
$OA=|r|\sqrt{a^2+b^2}$、
$OB=|r|^{-1}\sqrt{a^2+b^2}$，故
$\boxed{OA\cdot OB=a^2+b^2}$。
\end{proofbox}
\clearpage
\begin{proofbox}
\textbf{方法三｜常规解析几何（直线斜截式＋相切判别式）。}
\begin{center}
\begin{tikzpicture}[x=1.12cm,y=1.12cm,>=stealth,every node/.style={font=\small}]
\coordinate (O) at (0,0);
\coordinate (A) at (3.24,1.98);
\coordinate (B) at (1,-0.611111);
\coordinate (P) at (2.12,0.684444);
\fill[red!7] (O)--(A)--(B)--cycle;
\draw[gray!65,->] (-0.65,0)--(4.15,0) node[right] {$x$};
\draw[gray!65,->] (0,-1.35)--(0,2.6) node[above] {$y$};
\draw[blue!65!black,densely dashed] (-0.65,-0.3972)--(3.8,2.3222);
\draw[blue!65!black,densely dashed] (-0.65,0.3972)--(2.05,-1.2528);
\draw[teal!70!black,thick,domain=-0.9:1.22,samples=90,variable=\t]
 plot ({1.8/2*(exp(\t)+exp(-\t))},{1.1/2*(exp(\t)-exp(-\t))});
\draw[red!75!black,very thick] ($(B)!-0.15!(A)$)--($(A)!-0.12!(B)$);
\draw[orange!90!black,thick] (A)--(O)--(B);
\fill (O) circle(1.5pt) node[below left=5pt] {$O$};
\fill (A) circle(1.5pt) node[above left] {$A$};
\fill (B) circle(1.5pt) node[below=9pt] {$B$};
\fill (P) circle(1.5pt) node[left=4pt] {$P$};

\draw[gray,densely dotted] (A)--(3.24,0);
\draw[gray,densely dotted] (B)--(1,0);
\node[below=5pt] at (3.24,0) {$x_A$};
\node[above=5pt] at (1,0) {$x_B$};
\node[anchor=west] at (4.55,2.05) {\textcolor{red!75!black}{$y=kx+d$}};
\node[anchor=west] at (4.55,1.1) {$y_A=\frac ba x_A$};
\node[anchor=west] at (4.55,0.15) {$y_B=-\frac ba x_B$};
\node[anchor=west] at (4.55,-0.8) {$x_Ax_B=a^2$};
\end{tikzpicture}

\end{center}

先看非竖直切线 $y=kx+d$。图中投影给出两交点横坐标；代入双曲线，得到
\[
(b^2-a^2k^2)x^2-2a^2kdx-a^2(d^2+b^2)=0.
\]
有限切点的非竖直切线不平行于渐近线，二次项非零。相切时判别式为零，化简得
\[
\boxed{d^2=a^2k^2-b^2}\qquad
\left(|k|>\frac ba\right).
\]
与两渐近线相交，横坐标分别为
\[
x_A=-\frac{d}{k-b/a},\qquad
x_B=-\frac{d}{k+b/a},\qquad x_Ax_B=a^2.
\]
因为 $A=(x_A,(b/a)x_A)$、$B=(x_B,-(b/a)x_B)$，所以
\[
S_{\triangle AOB}=\frac ba|x_Ax_B|=ab.
\]
\textbf{竖直切线必须单独处理：} $x=\pm a$，相应两交点为
$(\pm a,\pm b)$、$(\pm a,\mp b)$，面积仍为 $ab$。
\textbf{优势：}题目先给一条未知切线 $y=kx+d$ 时可直接使用。

\medskip
\end{proofbox}
\clearpage
\begin{proofbox}
\textbf{方法四｜补充：渐近线坐标变换（揭示面积不变量）。}
\begin{center}
% In the transformed plane xi*eta=1, P'=(1.5,2/3).
\begin{tikzpicture}[x=1.25cm,y=1.25cm,>=stealth,every node/.style={font=\small}]
\coordinate (O) at (0,0);\coordinate (A) at (3,0);
\coordinate (B) at (0,1.333333);\coordinate (P) at (1.5,.666667);
\fill[red!7] (O)--(A)--(B)--cycle;
\draw[blue!65!black,->] (-.3,0)--(3.85,0) node[right] {$\xi$};
\draw[blue!65!black,->] (0,-.3)--(0,2.65) node[above] {$\eta$};
\draw[teal!70!black,thick,domain=.4:3.3,samples=90] plot (\x,{1/\x});
\draw[red!75!black,very thick] (-.25,1.444444)--(3.3,-.133333);
\draw[gray,densely dotted] (P)--(1.5,0) node[below] {$\xi_0$};
\draw[gray,densely dotted] (P)--(0,.666667) node[left] {$\eta_0$};
\fill (O) circle(1.5pt) node[below left] {$O'$};
\fill (A) circle(1.5pt) node[below] {$A'(2\xi_0,0)$};
\fill (B) circle(1.5pt) node[left=8pt] {$B'(0,2\eta_0)$};
\fill (P) circle(1.5pt) node[above right] {$P'(\xi_0,\eta_0)$};
\node[anchor=west] at (4.55,2.1) {$\xi\eta=1$};
\node[anchor=west] at (4.55,1.25) {$\xi_0\eta_0=1$};
\node[anchor=west] at (4.55,.4) {$S'=\frac12(2\xi_0)(2\eta_0)=2$};
\node[anchor=west] at (4.55,-.45) {图示右支，左支取绝对值};
\end{tikzpicture}

\end{center}

设 $\xi=x/a+y/b,\ \eta=x/a-y/b$，则
\[
H:\ \xi\eta=1,\qquad \ell_1:\eta=0,\quad\ell_2:\xi=0.
\]
切点为 $(\xi_0,\eta_0)$ 且 $\xi_0\eta_0=1$。
曲线 $\xi\eta=1$ 在该点的切线为
$\eta_0\xi+\xi_0\eta=2$，两轴截距为
$2\xi_0,2\eta_0$，所以新坐标中的三角形面积是
$\frac12|2\xi_0\cdot2\eta_0|=2$。
由变换行列式 $\bigl|\partial(\xi,\eta)/\partial(x,y)\bigr|=2/(ab)$，
原坐标的面积为 $2\cdot ab/2=ab$。
\textbf{优势：}理解``双曲线＋两渐近线''与 $xy=1$ 是同一结构；
用于后续研究一般仿射变换与面积比。高考考场优先方法一、二。

\end{proofbox}
